\chapter{Algèbre déployée et factorisation exacte du caractère}
\label{chap:factorizacion-split}
\index{algèbre déployée}\index{caractère!factorisation déployée}
\index{cône nul}\index{canaux $q_\pm$}

La signature formelle d'une extension et les canaux angulaires proviennent de branches déjà construites. Dans cette section, ils convergent vers une factorisation déterminantielle qui conserve séparément l'échelle radiale et l'orientation relative. Le résultat appartient au niveau algébrique : il n'identifie à lui seul ni les idéaux nuls à des photons physiques ni l'intérieur positif à une particule matérielle.

\section{Coordonnées angulaires conjuguées et réseau spectral bidirectionnel}
\label{sec:reticula-angular-split}

La chaîne causale précédente fixe
\[
 A_{\rm rad}=\frac{\pi A}{180},\qquad
 \Cstarrad=\frac{\pi\Cstar}{180},
\]
et, seulement après l'obtention de la coordonnée directe \(A\) et de l'unique coordonnée conjuguée \(C^*\),
\[
 q_+=\exp[-(A_{\rm rad}+\Cstarrad)],
 \qquad
 q_-=\exp[-(A_{\rm rad}-\Cstarrad)].
\]
Pour une signature angulaire \((n_A,s_C)\), introduisons les deux nombres d'enroulement
\[
 W_+=6n_A+s_C,
 \qquad
 W_-=6n_A-s_C.
\]

\begin{theorem}[Réseau angulaire d'indice douze]
\label{thm:reticula-two-way}
L'application
\[
 \mathcal W:\ZZ^2\to\ZZ^2,
 \qquad
 \binom{n_A}{s_C}\longmapsto
 \binom{W_+}{W_-}
 =\begin{pmatrix}6&1\\6&-1\end{pmatrix}
 \binom{n_A}{s_C}
\]
est injective et son image est d'indice douze. Ses facteurs invariants sont \((1,12)\), et
\[
 \operatorname{im}\mathcal W
 =\{(u,v)\in\ZZ^2:u+v\equiv0\pmod{12}\}.
\]
De plus,
\[
 \boxed{
 \exp\!\left(n_AA_{\rm rad}+\frac{s_C}{6}\Cstarrad\right)
 =\left(q_+^{-W_+}q_-^{-W_-}\right)^{1/12}.}
\]
\end{theorem}

\begin{proof}
La matrice a pour déterminant \(-12\). Le plus grand commun diviseur de ses coefficients étant égal à un, sa forme normale de Smith est \(\operatorname{diag}(1,12)\). Si \((u,v)=(6n+s,6n-s)\), alors \(u+v=12n\). Réciproquement, si \(u+v\) est un multiple de douze, alors
\[
 n=\frac{u+v}{12},
 \qquad
 s=\frac{u-v}{2}
\]
sont entiers et permettent de retrouver \((u,v)\). Enfin,
\[
 \frac{W_+(A_{\rm rad}+\Cstarrad)
       +W_-(A_{\rm rad}-\Cstarrad)}{12}
 =n_AA_{\rm rad}+\frac{s_C}{6}\Cstarrad,
\]
ce qui démontre l'identité par exponentiation.
\end{proof}

L'involution \(s_C\mapsto-s_C\) échange \(W_+\leftrightarrow W_-\) et \(q_+\leftrightarrow q_-\). La racine douzième correspond à l'indice entier de l'image de \(\mathcal W\), et non à un arrondi. La coordonnée \(C^*\) agit comme composante antisymétrique du réseau ; ce n'est ni une correction décimale d'une masse ni une sortie de \(\Gamma_{6\to8}\).

\section{Algèbre réelle déployée}
\label{sec:algebra-real-partida}

\begin{definition}[Algèbre réelle déployée]
Soit
\[
 \mathcal A_{\rm sp}
 :=\RR[\mathsf R]/(\mathsf R^2-1).
\]
La conjugaison et la norme sont définies respectivement par
\[
 \overline{a+b\mathsf R}=a-b\mathsf R,
 \qquad
 N_{\rm sp}(z)=z\overline z=a^2-b^2.
\]
\end{definition}

Les idempotents conjugués
\[
 P_\pm=\frac{1\pm\mathsf R}{2}
\]
satisfont
\[
 P_\pm^2=P_\pm,\qquad
 P_+P_-=0,\qquad
 P_++P_-=1,\qquad
 \overline{P_+}=P_-.
\]

\begin{proposition}[Décomposition spectrale déployée]
\label{prop:descomposicion-espectral-partida} Tout \(z\in\mathcal A_{\rm sp}\) admet une décomposition unique
\[
 z=r_+P_++r_-P_-,
 \qquad r_\pm\in\RR,
\]
et sa norme est
\[
 \boxed{N_{\rm sp}(z)=r_+r_-.}
\]
Chaque idéal principal \(\mathcal A_{\rm sp}P_\pm=\RR P_\pm\) est une droite nulle ; en particulier,
\[
 N_{\rm sp}(\lambda P_\pm)=0.
\]
\end{proposition}

\begin{proof}
Si \(z=a+b\mathsf R\), le choix \(r_+=a+b\), \(r_-=a-b\) fournit la décomposition, dont l'unicité découle de l'indépendance de \(P_+\) et \(P_-\). La conjugaison échange les deux coefficients : \(\overline z=r_-P_++r_+P_-\). L'orthogonalité des idempotents donne
\[
 z\overline z=(r_+r_-)(P_++P_-)=r_+r_-.
\]
L'énoncé relatif à chaque idéal s'obtient en annulant l'un des deux coefficients.
\end{proof}

\begin{definition}[Cône causal déployé]
\label{def:cono-causal-partido} Le cône positif, sa frontière nulle et son intérieur sont
\[
 \mathcal C_{\rm sp}^{+}
 =\{r_+P_++r_-P_-:r_+,r_-\geq0\},
\]
\[
 \partial_{\rm null}\mathcal C_{\rm sp}^{+}
 =\RR_{\geq0}P_+\cup\RR_{\geq0}P_-,
\qquad
 \operatorname{int}\mathcal C_{\rm sp}^{+}
 =\{r_+P_++r_-P_-:r_+r_->0\}.
\]
\end{definition}

La frontière possède une seule direction active et une norme nulle. L'intérieur exige deux directions simultanées et une norme positive. Cette stratification est exacte ; sa réalisation physique exige des applications qui préservent la norme, le transport et la dynamique concernés.

\section{Transporteur relatif et conditions de réalisation représentationnelle}
\label{sec:transportador-relativo-split-md}

Le bloc historique \(B_0=7I+2\mathsf R\) a pour spectre \(\{9,5\}\), et le bloc angulaire \(B_1=AI+\Cstar\mathsf R\) a pour spectre \(\{A+\Cstar,A-\Cstar\}\). Le \cref{thm:no-go-entrelazador-bloques-split} démontre qu'avec les valeurs internes en vigueur, tout entrelaceur linéaire supposé \(\Phi B_0=B_1\Phi\) est nul. L'objet actif est le transporteur multiplicatif relatif déjà construit, et non une équivalence entre représentations de spectres disjoints.

\section{Factorisation déployée du caractère angulaire}
\label{sec:factorizacion-split-angular}

Soit
\[
 \sigma=\frac{s_C}{6},\qquad
 w_\pm=\frac{n_A\pm\sigma}{2}=\frac{W_\pm}{12},
 \qquad
 r_\pm=q_\pm^{-w_\pm},
\]
et définissons
\[
 Z_v=r_+P_++r_-P_-.
\]

\begin{theorem}[Factorisation déployée du caractère angulaire]
\label{thm:split-character} Pour toute signature angulaire \(v=(n_A,s_C)\in\ZZ^2\),
\[
 \boxed{
 N_{\rm sp}(Z_v)
 =\det Z_v
 =\exp\!\left(
 n_AA_{\rm rad}+\frac{s_C}{6}\Cstarrad
 \right).}
\]
\end{theorem}

\begin{proof}
Par définition des canaux,
\[
 \log r_+
 =\frac{n_A+\sigma}{2}(A_{\rm rad}+\Cstarrad),
 \qquad
 \log r_-
 =\frac{n_A-\sigma}{2}(A_{\rm rad}-\Cstarrad).
\]
En sommant,
\[
 \log(r_+r_-)
 =n_AA_{\rm rad}+\sigma\Cstarrad.
\]
La \cref{prop:descomposicion-espectral-partida} donne \(N_{\rm sp}(Z_v)=r_+r_-\), et la représentation diagonale dans la base \((P_+,P_-)\) donne la même expression pour le déterminant.
\end{proof}

La composante que le déterminant ne détecte pas conserve l'orientation relative :
\[
 \frac12\log\frac{r_+}{r_-}
 =\frac12\left(
 n_A\Cstarrad+\frac{s_C}{6}A_{\rm rad}
 \right),
\]
\[
 Z_v=
 \exp\!\left[
 \frac12\left(n_AA_{\rm rad}+\frac{s_C}{6}\Cstarrad\right)
 \right]
 \exp\!\left[
 \frac12\left(n_A\Cstarrad+\frac{s_C}{6}A_{\rm rad}\right)
 \mathsf R
 \right].
\]
Le premier facteur fixe l'échelle radiale ; le second enregistre l'orientation relative. Ce sont deux lectures coordonnées d'un même élément.

\section{Incorporation centrale des tours}
\label{sec:factorizacion-split-torres}

Une fois \(\Dfour\), \(s_{120}\) et \(\zCorr=\zeta_{270}^{\rm corr}\) construits globalement, soit
\[
 t(v)=
 k_{\Dfour}\Dfour
 +\nu_{120}s_{120}
 +\nu_{270}\zCorr,
 \qquad
 \widetilde Z_v=e^{t(v)/2}Z_v.
\]

\begin{corollary}[Caractère complet comme norme déployée]
\label{cor:split-character-torres} Les blocs globaux précédents étant fixés,
\[
 \boxed{
 N_{\rm sp}(\widetilde Z_v)
 =\exp\!\left(
 n_AA_{\rm rad}+\frac{s_C}{6}\Cstarrad
 +k_{\Dfour}\Dfour+\nu_{120}s_{120}
 +\nu_{270}\zCorr
 \right).}
\]
\end{corollary}

\begin{proof}
Dans la représentation bidimensionnelle, la multiplication par le scalaire central \(e^{t(v)/2}\) multiplie le déterminant par \(e^{t(v)}\). On applique le \cref{thm:split-character}.
\end{proof}

Le corollaire construit la réalisation déployée déclarée du caractère actif ; il ne classe pas toutes les réalisations possibles et ne transforme pas les tours en preuves indépendantes. Ses hypothèses sont les canaux et les blocs globaux déjà fixés, et son contrôle négatif est immédiat : omettre le facteur central \(1/2\) double l'exposant de tour ; activer simultanément les deux idempotents avec des coefficients positifs fait sortir de la frontière nulle.
